\subsection{有理线性形式}

设 V 是 K 上的一个右向量空间，带有一个 \(K^{\prime}\) -结构 \(V^{\prime}\) 。由于 \(K_{d}^{\prime}\) 是一个 \(K^{\prime}\) -结构（在右向量 K-空间 \(K_{d}\) 上的），我们可以把线性形式 \(x^{*} \in V^{*}\) （在 \(K^{\prime}\) 上有理的）定义为 V 到 \(K_{d}\) 中的线性映射（在 \(K^{\prime}\) 上有理的，相对于

上的 K'-结构，即 V 上的和 \(K_{d}\) 上的）。根据第 3 号命题 3，这些线性形式的集合 \(R'\) 是对偶空间 \(V'^*\) （ \(V'\) 的对偶）在复合映射

\[
\mathrm{V} ^ {\prime *} \xrightarrow {\phi} \mathrm{K} \otimes_ {\mathrm{K} ^ {\prime}} \mathrm{V} ^ {\prime *} \xrightarrow {\upsilon} \mathrm{V} ^ {*}\tag{1}
\]

下的像，其中 \(\phi(x^{\prime*}) = 1 \otimes x^{\prime*}\) ，而 \(\upsilon(\xi \otimes x^{\prime*})\) 是 V 上的线性形式 \(y^{*}\) ，使得 \(y^{*}(x^{\prime}) = \xi\langle x^{\prime*}, x^{\prime} \rangle\) 对所有 \(x^{\prime} \in V^{\prime}\) 成立（§ 5，第 4 号）。我们知道这个映射是单射（§ 7，第 9 号，命题 20 和 19），且显然 \(R^{\prime}\) 是一个左向量子 \(K^{\prime}\) -空间（ \(V^{*}\) 的）；此外， \(R^{\prime}\) 的每个在 \(K^{\prime}\) 上自由的子集在 K 上也是自由的。但一般说来， \(R^{\prime}\) 不一定在 K 上生成 \(V^{*}\) ，因而并不定义一个 \(K^{\prime}\) -结构于 \(V^{*}\) 之上（习题 2）。然而，若 V 在 K 上是有限维 \(n\) 的，则 \(V^{\prime*}\) 是维数 \(n\) 的（在 \(K^{\prime}\) 上），于是 \(R^{\prime}\) 典范地定义一个 \(K^{\prime}\) -结构于 \(V^{*}\) 之上。

命题 5. 设 V 是 K 上的右向量空间，V' 是 V 上的一个 K'-结构，W 是 V 的一个向量子 K-空间。要使 W 在 K' 上有理，必要且充分的条件是：存在一个由在 K' 上有理的线性形式组成的集合 H ⊂ V*，使得 W 是 H 在 V 中的正交补（§ 2，第 4 号）。

设 H 是 V* 的一个子集，其元素是在 K' 上有理的线性形式。对所有 \(x^{*} \in H\) ， \(x^{*}\) 的核是 V 的一个在 K' 上有理的向量子 K-空间（第 3 号，命题 3 的推论 2）；因此这些核的交也是 V 的一个在 K' 上有理的向量子 K-空间（第 2 号，命题 2 的推论 1）。

反之，设 W 是 V 的一个在 \(K'\) 上有理的向量子 K-空间，于是 W 等同于 \(W' \otimes_{K'} K\) ，其中 \(W' = W \cap V'\) （第 2 号，命题 2）。要使线性形式 \(x'^* \in V'^*\) 在 \(W'\) 上为零，必要且充分的条件是它在 (1) 下对应的线性形式 \(x^* \in V^*\) 在 W 上为零，因为根据第 3 号命题 3 的推论 1，

\[
\operatorname{Ker} (x ^ {*}) = \left(\operatorname{Ker} (x ^ {\prime *})\right) \otimes_ {\mathbb {K} ^ {\prime}} \mathrm{K} \quad \text { and } \quad \operatorname{Ker} (x ^ {* ^ {\prime}}) = (\operatorname{Ker} (x ^ {*})) \cap \mathrm{V} ^ {\prime}.
\]

设 \(H'\) 是 \(W'\) 在 \(V'^*\) 中的正交补；我们知道（§ 7，第 5 号，定理 7）， \(W'\) 是 \(H'\) 在 \(V'\) 中的正交补；若 \(H\) 是 \(H'\) 在映射 (1) 下在 \(V*\) 中的像，则由上述可知， \(W\) 是 \(H\) 在 \(V\) 中的正交补，这里用到第 7 号命题 14 的推论。

